% OpenStax Calculus Volume 2 -- Section 6.3 Exercises: Taylor and Maclaurin Series
% Exercises reproduced from OpenStax, licensed under CC BY 4.0.
% Source: https://openstax.org/books/calculus-volume-2/pages/6-3-taylor-and-maclaurin-series
\documentclass{exercises}
\title{OpenStax Calculus Volume 2}
\subtitle{Section 6.3 Exercises: Taylor and Maclaurin Series}
\sourceurl{https://openstax.org/books/calculus-volume-2/pages/6-3-taylor-and-maclaurin-series}
\author{}
\date{}

\begin{document}
\maketitle


In the following exercises, find the Taylor polynomials of degree two approximating the given function centered at the given point.

\begin{enumerate}[start=1]
\item $f(x)=1+x+x^{2}$ at $a=1$
\item $f(x)=1+x+x^{2}$ at $a=-1$
\item $f(x)=\cos(2x)$ at $a=\pi$
\item $f(x)=\sin(2x)$ at $a=\frac{\pi}{2}$
\item $f(x)=\sqrt{x}$ at $a=4$
\item $f(x)=\ln x$ at $a=1$
\item $f(x)=\frac{1}{x}$ at $a=1$
\item $f(x)=e^{x}$ at $a=1$
\end{enumerate}

In the following exercises, verify that the given choice of $n$ in the remainder estimate $|R_{n}|\le \frac{M}{(n+1)!}(x-a)^{n+1}$, where $M$ is the maximum value of $|f^{(n+1)}(z)|$ on the interval between $a$ and the indicated point, yields $|R_{n}|\le \frac{1}{1000}$. Find the value of the Taylor polynomial $p_{n}$ of $f$ at the indicated point. If $|R_{n}|$ is not less than $1/1000$, determine what it is.

\begin{enumerate}[resume]
\item \techrequired $\sqrt{10};\ a=9,\ n=3$
\item \techrequired $28^{1/3};\ a=27,\ n=1$
\item \techrequired $\sin(6);\ a=2\pi,\ n=5$
\item \techrequired $e^{2};\ a=0,\ n=9$
\item \techrequired $\cos\left(\frac{\pi}{5}\right);\ a=0,\ n=4$
\item \techrequired $\ln(2);\ a=1,\ n=1000$
\item Integrate the approximation $\sin t\approx t-\frac{t^{3}}{6}+\frac{t^{5}}{120}-\frac{t^{7}}{5040}$ evaluated at $\pi t$ to approximate $\int_{0}^{1}\frac{\sin(\pi t)}{\pi t}\,dt$.
\item Integrate the approximation $e^{x}\approx1+x+\frac{x^{2}}{2}+\cdots+\frac{x^{6}}{720}$ evaluated at $-x^{2}$ to approximate $\int_{0}^{1}e^{-x^{2}}\,dx$.
\end{enumerate}

In the following exercises, find the smallest value of $n$ such that the remainder estimate $|R_{n}|\le \frac{M}{(n+1)!}(x-a)^{n+1}$, where $M$ is the maximum value of $|f^{(n+1)}(z)|$ on the interval between $a$ and the indicated point, yields $|R_{n}|\le \frac{1}{1000}$ on the indicated interval.

\begin{enumerate}[resume]
\item $f(x)=\sin x$ on $[-\pi,\pi]$, $a=0$
\item $f(x)=\cos x$ on $[-\frac{\pi}{2},\frac{\pi}{2}]$, $a=0$
\item $f(x)=e^{-2x}$ on $[-1,1]$, $a=0$
\item $f(x)=e^{-x}$ on $[-3,3]$, $a=0$
\end{enumerate}

In the following exercises, the maximum of the right-hand side of the remainder estimate $|R_{1}|\le \frac{\max |f''(z)|}{2}R^{2}$ on $[a-R,a+R]$ occurs at $a$ or $a\pm R$. Estimate the maximum value of $R$ such that $\frac{\max |f''(z)|}{2}R^{2}\le 0.1$ on $[a-R,a+R]$ by plotting this maximum as a function of $R$.

\begin{enumerate}[resume]
\item \techrequired $e^{x}$ approximated by $1+x$, $a=0$
\item \techrequired $\sin x$ approximated by $x$, $a=0$
\item \techrequired $\ln x$ approximated by $x-1$, $a=1$
\item \techrequired $\cos x$ approximated by $1$, $a=0$
\end{enumerate}

In the following exercises, find the Taylor series of the given function centered at the indicated point.

\begin{enumerate}[resume]
\item $x^{4}$ at $a=-1$
\item $1+x+x^{2}+x^{3}$ at $a=-1$
\item $\sin x$ at $a=\pi$
\item $\cos x$ at $a=2\pi$
\item $\sin x$ at $x=\frac{\pi}{2}$
\item $\cos x$ at $x=\frac{\pi}{2}$
\item $e^{x}$ at $a=-1$
\item $e^{x}$ at $a=1$
\item $\frac{1}{(x-1)^{2}}$ at $a=0$ (Hint: Differentiate $\frac{1}{1-x}$.)
\item $\frac{1}{(x-1)^{3}}$ at $a=0$
\item $F(x)=\int_{0}^{x}\cos(\sqrt{t})\,dt$; $f(t)=\sum_{n=0}^{\infty}(-1)^{n}\frac{t^{n}}{(2n)!}$ at $a=0$ (Note: $f$ is the Taylor series of $\cos(\sqrt{t})$.)
\end{enumerate}

In the following exercises, compute the Taylor series of each function around $x=1$.

\begin{enumerate}[resume]
\item $f(x)=2-x$
\item $f(x)=x^{3}$
\item $f(x)={(x-2)}^{2}$
\item $f(x)=\ln x$
\item $f(x)=\frac{1}{x}$
\item $f(x)=\frac{1}{2x-x^{2}}$
\item $f(x)=\frac{x}{4x-2x^{2}-1}$
\item $f(x)=e^{-x}$
\item $f(x)=e^{2x}$
\end{enumerate}

\techrequired In the following exercises, identify the value of $x$ such that the given series $\sum_{n=0}^{\infty}a_{n}$ is the value of the Maclaurin series of $f$ at $x$. Approximate $f(x)$ using $S_{10}=\sum_{n=0}^{10}a_{n}$.

\begin{enumerate}[resume]
\item $\sum_{n=0}^{\infty}\frac{1}{n!}$
\item $\sum_{n=0}^{\infty}\frac{2^{n}}{n!}$
\item $\sum_{n=0}^{\infty}\frac{{(-1)}^{n}{(2\pi)}^{2n}}{(2n)!}$
\item $\sum_{n=0}^{\infty}\frac{{(-1)}^{n}{(2\pi)}^{2n+1}}{(2n+1)!}$
\end{enumerate}

The following exercises make use of the functions $S_{5}(x)=x-\frac{x^{3}}{6}+\frac{x^{5}}{120}$ and $C_{4}(x)=1-\frac{x^{2}}{2}+\frac{x^{4}}{24}$ on $[-\pi,\pi]$.

\begin{enumerate}[resume]
\item \techrequired Plot $\sin^{2}x-{(S_{5}(x))}^{2}$ on $[-\pi,\pi]$. Compare the maximum difference with the square of the Taylor remainder estimate for $\sin x$.
\item \techrequired Plot $\cos^{2}x-{(C_{4}(x))}^{2}$ on $[-\pi,\pi]$. Compare the maximum difference with the square of the Taylor remainder estimate for $\cos x$.
\item \techrequired Plot $|2S_{5}(x)C_{4}(x)-\sin(2x)|$ on $[-\pi,\pi]$.
\item \techrequired Compare $\frac{S_{5}(x)}{C_{4}(x)}$ on $[-1,1]$ to $\tan x$. Compare this with the Taylor remainder estimate for the approximation of $\tan x$ by $x+\frac{x^{3}}{3}+\frac{2x^{5}}{15}$.
\item \techrequired Plot $e^{x}-e_{4}(x)$ where $e_{4}(x)=1+x+\frac{x^{2}}{2}+\frac{x^{3}}{6}+\frac{x^{4}}{24}$ on $[0,2]$. Compare the maximum error with the Taylor remainder estimate.
\item (Taylor approximations and root finding.) Recall that Newton's method $x_{n+1}=x_{n}-\frac{f(x_{n})}{f'(x_{n})}$ approximates solutions of $f(x)=0$ near the input $x_{0}$.
\begin{enumerate}[label=(\alph*)]
  \item If $f$ and $g$ are inverse functions, explain why a solution of $g(x)=a$ is the value $f(a)$ of $f$.
  \item Let $p_{N}(x)$ be the degree-$N$ Maclaurin polynomial of $e^{x}$. Use Newton's method to approximate solutions of $p_{N}(x)-2=0$ for $N=4,5,6$.
  \item Explain why the approximate roots of $p_{N}(x)-2=0$ are approximate values of $\ln (2)$.
\end{enumerate}
\end{enumerate}

In the following exercises, use the fact that if $q(x)=\sum_{n=0}^{\infty}a_{n}(x-c)^{n}$ converges in an interval containing $c$, then $\lim_{x\to c}q(x)=a_{0}$ to evaluate each limit using Taylor series.

\begin{enumerate}[resume]
\item $\lim_{x\to0}\frac{\cos x-1}{x^{2}}$
\item $\lim_{x\to0}\frac{\ln (1-x^{2})}{x^{2}}$
\item $\lim_{x\to0}\frac{e^{x^{2}}-x^{2}-1}{x^{4}}$
\item $\lim_{x\to0^{+}}\frac{\cos(\sqrt{x})-1}{2x}$
\end{enumerate}

\end{document}
